\documentclass[12pt,a4paper]{article}

%----------------------------------------------------------------------
% Page layout
%----------------------------------------------------------------------
\usepackage[margin=1in]{geometry}
\usepackage{setspace}
\doublespacing
%\singlespacing % use while the paper is still an outline

%----------------------------------------------------------------------
% Fonts and encoding
%----------------------------------------------------------------------
\usepackage[T1]{fontenc}
\usepackage[utf8]{inputenc}
\usepackage{newtxtext,newtxmath}
\usepackage{microtype}

%----------------------------------------------------------------------
% Maths
%----------------------------------------------------------------------
\usepackage{amsmath,amssymb}

%----------------------------------------------------------------------
% Tables and figures
%----------------------------------------------------------------------
\usepackage{graphicx}
\usepackage{booktabs}
\usepackage{tabularx}
\usepackage{longtable}
\usepackage{threeparttable}
\usepackage{rotating}
\usepackage[labelfont=bf,font=small]{caption}
\usepackage{float}

%----------------------------------------------------------------------
% Lists
%----------------------------------------------------------------------
\usepackage{enumitem}

%----------------------------------------------------------------------
% Citations
%----------------------------------------------------------------------
\usepackage[round,authoryear]{natbib}

%----------------------------------------------------------------------
% Miscellaneous
%----------------------------------------------------------------------
%\usepackage[nolists,tablesfirst]{endfloat} % uncomment to move floats to the end

%----------------------------------------------------------------------
% hyperref last
%----------------------------------------------------------------------
\usepackage[hidelinks,breaklinks=true]{hyperref}

%----------------------------------------------------------------------
% Front matter
%----------------------------------------------------------------------
\title{Rupture or Continuity: Assessing the 2025 Structural Break in World Politics}
\author{Max Gallop and Shahryar Minhas} % remove for blind review
\date{\today}

\begin{document}

\maketitle

\begin{abstract}
\noindent For almost 40 years, international politics has been shaped and defined by the liberal, Western-led order. However, since the beginning of the second Trump administration, it feels as if we are experiencing what Canadian Prime Minister Mark Carney called a ``rupture'' in world politics. In this paper we seek to assess whether 2025 represented a genuine rupture, or the intensification and acceleration of long-running trends. We do this by using a novel network model which is particularly designed to allow for ``snap-shift'' dynamics in a latent space, and applying this model to a variety of measures of international economic and diplomatic cooperation. We also test how well models trained on the 1989--2024 period do at explaining state behaviour in 2025 and 2026, which allows us to assess the extent to which data-driven guidance needs to be re-evaluated in our rapidly changing landscape.
\end{abstract}

\bigskip
\noindent \textbf{Keywords:} international order; alignment; network models; latent space; structural break

\clearpage

\section{Introduction}
\begin{itemize}
\item Hook: craziness of the first 18 months of the second Trump admin, especially in IR. Pair with increasing Russia/China influence and cooperation.
\item Carney rupture quote.
\item General tendency to believe we're living in unprecedented times.
\end{itemize}

\subsection{Our Goal in this paper}

Principled assessment of whether there has been a rupture in international alignment, using network approaches.

\begin{enumerate}
\item Subquestion 1: If there has been a rupture, is it universal or is there variance across forms of cooperation?
\item Subquestion 2: If there has been a rupture, what does the new world look like?
\item Subquestion 3: What happens if we assume there hasn't been a rupture?
\end{enumerate}

\section{Literature}
\begin{itemize}
\item Old school alignment stuff -- balancing/bandwagoning, power transition theory and international orders, hierarchy stuff.
\item More recent state preferences/alignment quantitative work.
\item Recent work about global south and Russia/China.
\end{itemize}

\section{Methodology}
\begin{itemize}
\item Why network models generally?
\item Discussion of AME models as a class.
\item How time has traditionally been handled.
\item Discussion of snap/shift stuff in this paper.
\end{itemize}

\section{Data}
\begin{itemize}
\item UN Voting
\item Trade
\item Sanctions
\item Event data?
\end{itemize}

\section{Results}

\subsection{What evidence is there of a rupture?}

\subsection{SQ1}

Comparison of snap-shift effects of different data sources.

\subsection{SQ2}

More detail on behaviour of non-US actors, specific focus on Ch/Rus and EU.

\subsection{SQ3}

Comparison to network model with non snap-shift. Or train model until 2024, then use to predict?

\section{Discussion}

The old world is dying, the new world is yearning to be born; now is the time of monsters. 

\clearpage
\bibliographystyle{apsr}
\bibliography{references}

\end{document}
